\graphicspath{{figures/intro}}
\section{Introduction}  \label{sec:intro}
\par At the center of most galaxies resides a super massive black hole (SMBH) with masses typically exceeding $10^6M_\odot$. When these SMBHs rapidly accrete surrounding gas and dust, they are classified as active galactic nuclei (AGN). Close to the center of these AGN, fueled primarily by the gravitational potential energy of the infalling material, a relatively small disk (the accretion disk) of super heated plasma emits primarily in the X-ray, ultraviolet (UV), and optical portions of the electromagnetic (EM) spectrum. These AGN have luminosities typically between $10^{40}$ ergs•s-1 and $10^{47}$ erg•s-1 and can, in some cases, exceed the luminosity of the host galaxy. Surrounding this accretion disk is a much larger structure of dust, which, when heated by the accretion disk, emits primarily in the infrared (IR). Understanding the structure and geometry of this dust provides insights into the growth and evolution of SMBHs and how AGN interact more broadly with their host galaxies allowing for a greater understanding of the processes governing galactic evolution. \zach{Need to revise this to be a little more specific. Definitely glossing over a lot of stuff that should be cited/quantified.}
\par However, due to the milliarcsecond spatial resolution required to directly resolve the inner radius of the dust distributions in even the closest AGN, any attempts to understand the hottest portion of the dust currently require some form of modeling. \zach{Citeeee} Multiple successful attempts have been made to resolve the inner radius of several close AGN utilzing IR interferometry, but this method does not extend well to AGN of higher redshift. {\bf [Cited in Almeyda et al. 2017: Swain et al. 2003; Kishimoto et al. 2009; Pott et al. 2010; Kishimoto et al. 2011; Weigelt et al. 2012], [Cited in Almeyda et al. 2020: e.g., Kishimoto et al. 2011; Burtscher et al. 2013]} \zach{This needs to be expanded to actually say what these people did}
\par Instead, an alternative approach, based on reverberation mapping, uses both optical and IR wavelengths to bypass the need for high resolution imaging by utilizing the time domain. \zach{Instead of just jumping into saying this method exists, maybe cite a time / the times when this was first used?} The driving principle of dust reverberation mapping is that the IR continuum response is delayed relative to the optical continuum emission since the optical light from the accretion disk must first travel to the dust where it is absorbed and then re-emitted in the IR along the observer’s path. \zach{Theory of reverberation mapping says that 
Theory of reverberation mapping
how do we implement it observationally
how do we extend it to modeling}
\par Observational reverberation mapping is commonly used to constrain the characteristic size of the torus { \bf [refs]} by determining the average[?] delay between the optical and IR light curves, then converting this time delay into a characteristic radius. In an attempt to gain additional information about the structure and geometry of the circumnuclear dust, we also consider that the total IR response at any given time is a summation of partial responses to the optical emission at all time delays. This means that the optical emission and IR response will not be perfectly correlated. We can instead model the IR response light curve as a convolution of the optical emission with a “transfer function”. This transfer function then would contain information about the structure and geometry of the dust.  

% Extending this principle, we exploit the fact that dust at different radii and at different angles relative to the central source will respond at different times.

% consider that dust at 
% different radii and at different
% angle between the cloud and the observer and the angle between cloud and the cenral source. 
In an attempt to gain more information about the structure, 
However, since the observed light curve is a convolution with the transfer function
transfer function contains information about the structure and geometry of the torus
we can use the transfer function to learn more about the 

Start the other way: we know that from reverberation mapping - transfer function - transfer function contains information - extract that information using modeling 

\par The forward modeling code TORMAC (TOrus Reverberation MApping Code) has previously been employed to simulate the multi-wavelength response from a circumnuclear dust distribution to a delta function “pulse”, emulating that dust configuration’s “transfer function” (Almeyda et al. 2017, 2020). However, it is also capable of generating IR response light curves to observed input optical light curves. \zach{I think you reminded me a while ago that TORMAC has already done this, need to rephrase} In this paper, we utilize TORMAC to extract additional information about the dust structure of a well-studied AGN, NGC 3783.
\par IR interferometric observations of NGC 3783 have revealed the presence of a hot, compact, disk-like structure, commonly referred to as the “dusty torus” and a cooler, resolved, extended component in the polar regions (GRAVITY Collaboration 2021 {\bf [among others]}). A  significant fraction of the IR emission from NGC 3783 comes from this polar dust structure, which we refer to as the “polar bi-cone". Both IR interferometry and SED model fitting { \bf [reffssss]} have been used to gain information about NGC 3783’s circumnuclear dust, but this paper serves as the first use of light curves to extract additional information about the circumnuclear dust.

% \begin{itemize}
%     % Which studies have tried? What were their results?
%     % \item Our approach is based on light echo modeling / reverberation mapping, which uses both optical and infrared (IR) wavelengths to bypass the need for high resolution imaging by utilizing the time domain. 
%     % \item A driving principal of light echo modeling is that the IR response is delayed since optical light from the accretion disk first travels to the dust before it is re-emitted in IR on the observer's path. 
%     % \item While observational reverberation mapping is commonly used to constrain the inner size of the torus, we can extend this principle by modeling the IR response over time as a convolution of the input optical data with a transfer function that depends on the dust’s structure and geometry. 
%     % \item This transfer function is difficult to determine from the observed data, so I utilize a python code TORMAC \citep{almeyda_modeling_2017,almeyda_modeling_2020} to attempt to extract additional information about the dust structure of a well studied AGN, NGC 3783. 
%     % Obviously need to provide details on what makes NGC 3783 well studied, what information do we already know
%     % \item NGC 3783 has had some of its torus characteristics constrained through IR interferometry, revealing an inclination that is close to being face-on to the observer \citep{gravity_collaboration_geometric_2021} and a maximum mid-IR radial extent of the dust \citep{as}(Asmus et al., 2014).
%     % \item AGN have a central disk of superheated gas called the accretion disk, and a dust structure which surrounds the accretion disk, separated into a torus (a donut shaped disk) and a polar bi-cone (a hollow cone at the north and south poles).
%     % % Add in a brief discussion of the evidence for polar wind
%     % \item Observations have shown that a significant fraction of the IR emission from NGC 3783 comes from a resolved, extended polar dust structure which is irreconcilable with torus only models. Recently, significant work has been done using {\bf SED?} model fitting to constrain the structure and geometry of this polar outflow. For NGC 3783, kinematic simulations agree on an extended (relative to the size of the torus) biconical outflow with an inclination to the observer LOS of 60-75°, an angular width of about 10°, and an opening angle of 30-45° \citep{fischer_determining_2013,muller-sanchez_outflows_2011}. 
%     % \item \cite{fischer_determining_2013} uses kinematics simulations for NGC3783 which fit well with an extended biconical outflow with $i_{POL} \approx 75^\circ$ ($i_{TOR} \approx 35^\circ$), $\sigma_{POL} \approx 10^{\circ}$, and $open \approx 50^{\circ}$. These findings are also supported by \cite{muller-sanchez_outflows_2011} which also suggests an extended biconical outflow with $i_{POL} \approx 60^\circ$ ($i_{TOR} \approx 38^{\circ}$) $\sigma_{POL} \approx 7^\circ$, and $open \approx 30^\circ$
%     % \item This is the first time a time domain radiative transfer model is being used to fit simulated to observed responses for AGN, and our method will be applied to many other reverberation mapped AGN.
% \end{itemize}